\documentclass[10pt]{article}

\usepackage[margin=1in]{geometry}
\usepackage{amsmath,amssymb,amsthm}
\usepackage{booktabs}
\usepackage{graphicx}
\usepackage{hyperref}
\usepackage{xcolor}

\title{Slot-Addressed Coupled Write-Forget Affine Transport for Recurrent State Tracking}
\author{Anonymous Authors}
\date{}

\newcommand{\mogt}{MOGT}
\newcommand{\R}{\mathbb{R}}

\begin{document}
\maketitle

\begin{abstract}
Transformers remain the dominant backbone for language modeling, but attention
couples sequence modeling to explicit pairwise token interactions. This paper
studies a non-attention alternative based on matrix-valued affine transport:
\[
    H_t = U_t H_{t-1} + V_t.
\]
The transition/value pair admits an associative composition law, enabling
parallel scan implementations while preserving recurrent state semantics. Our
initial language-modeling experiments do not support the broad claim that this
operator replaces Transformers. Instead, controlled synthetic experiments reveal
a sharper result: ungated skew-symmetric transport is poorly matched to
overwrite-style memory, while token-dependent gates and identity-like transport
make recurrent state tracking learnable. On a single-slot last-value task
trained at context 128, gated \mogt{} reaches 69.71\% $\pm$ 18.69\% accuracy at
context 1024 across three seeds, compared with 10.51\% $\pm$ 1.24\% for a
matched RoPE Transformer. However, a NoPE Transformer is a much stronger
position-robust baseline, reaching 68.64\% $\pm$ 15.51\% at the same length
under the matched 2048-example evaluation protocol.
At train context 512 with dense state supervision, an identity-transport,
dual-gated \mogt{} learns the task across all three seeds, but still trails NoPE
Transformer at far extrapolation. We then test a coupled write-forget gate,
which uses the same token-dependent scalar to write a new value and erase the
old state. This variant reaches 100.00\% $\pm$ 0.00\% accuracy from context 512
through 8192 after training only at context 512, compared with 78.65\% $\pm$
12.63\% for the current standard-schema NoPE Transformer rerun at 8192.
Finally, we add a slot-addressed gate
input that sees the current token, previous token, and tracked-slot prefix. On
a tracked 2-slot task, this reaches 86.46\% $\pm$ 3.25\% at context 4096 across
three seeds, compared with 44.79\% $\pm$ 9.42\% for NoPE Transformer. On a
tracked 4-slot task it reaches 98.96\% $\pm$ 0.90\% at context 4096, compared
with 24.48\% $\pm$ 7.05\% for NoPE Transformer. These results do not establish
a general Transformer replacement. They do, however, extend beyond dense
supervision: with final-query-only supervision and a 2-to-4-slot curriculum,
the 4-slot variant reaches 94.27\% $\pm$ 2.39\% at context 4096, compared with
21.35\% $\pm$ 6.31\% for the current standard-schema NoPE Transformer rerun
and 19.79\% $\pm$ 10.97\% for parameter-matched HF-Mamba d192; a 6-slot
version reaches 96.88\% $\pm$ 2.71\% versus 21.88\% $\pm$ 8.27\%, and an
8-slot version
reaches 85.42\% $\pm$ 11.93\%, versus 10.94\% $\pm$ 2.71\% for NoPE
Transformer, 18.23\% $\pm$ 1.80\% for HF Mamba at $d=128$, and 16.15\% $\pm$
6.31\% for a parameter-matched HF Mamba at $d=192$. The isolated mechanism is
concrete: affine transport needs write-aligned forgetting and task-relevant
addressing signals to preserve routed recurrent state.
\end{abstract}

\section{Introduction}

Self-attention made it possible to train large sequence models with strong
parallelism and flexible token-to-token interactions \cite{vaswani2017attention}.
Its success, however, comes with an explicit interaction pattern over the
sequence. This has motivated a long line of recurrent, linear-attention, and
state-space alternatives that seek better scaling at long context while
retaining parallel training \cite{katharopoulos2020transformers,gu2021s4,
sun2023retnet,peng2023rwkv,gu2023mamba,koo2023gateloop}.

This project began with an ambitious question: can a matrix-valued transport
operator replace attention? The current evidence says: not yet. On a
short-budget WikiText-103 protocol, a matched scratch Transformer is stronger
than the current \mogt{} language model. Rather than hiding that negative
result, we use it to narrow the claim. The useful contribution is not a broad
Transformer overthrow; it is a mechanistic study of when affine recurrent
transport works, when it fails, and what gates it needs.

The operator we study evolves a rank-valued state $H_t \in \R^{r \times c}$ by
a matrix transition $U_t \in \R^{r \times r}$ and value update
$V_t \in \R^{r \times c}$:
\begin{equation}
    H_t = U_t H_{t-1} + V_t.
    \label{eq:affine}
\end{equation}
This recurrence is attractive because transition/value pairs compose
associatively. It can therefore be trained with scan-style implementations
rather than a strictly sequential loop.

Our main empirical findings are:
\begin{itemize}
    \item Ungated \mogt{} and Cayley/Magnus-style transport stay near chance on
    overwrite memory tasks.
    \item Adding token-dependent forget/write gates makes the operator
    learnable on single-slot state tracking.
    \item Identity transport is much better than learned orthogonal-like
    transport for overwrite-style memory.
    \item On train-context-512 dense supervision, coupling the write gate to the
    forget gate yields 100\% accuracy through context 8192 across three seeds,
    outperforming the NoPE Transformer baseline on this synthetic task.
    \item Prefix-conditioned slot-addressed gates turn multi-slot tracking from
    a weak signal into a strong result on 2-slot and 4-slot probes.
    \item A slot-count curriculum makes 4-slot, 6-slot, and 8-slot final-query-only training
    reliable across three seeds, separating the operator's routing bias from
    dense state supervision alone.
\end{itemize}

\section{Affine Transport}

\subsection{Associative Composition}

Define an affine transition as a pair $(U,V)$ acting on state $H$ by
\[
    (U,V) \cdot H = UH + V.
\]
Two consecutive transitions compose as
\begin{equation}
    (U_2,V_2) \circ (U_1,V_1)
    =
    (U_2 U_1,\; U_2 V_1 + V_2).
    \label{eq:compose}
\end{equation}
This composition is associative:
\[
    ((U_3,V_3)\circ(U_2,V_2))\circ(U_1,V_1)
    =
    (U_3,V_3)\circ((U_2,V_2)\circ(U_1,V_1)).
\]
Equation~\ref{eq:compose} is the systems hook: prefix states can be computed by
a parallel scan over affine pairs.

\subsection{MOGT Block}

The \mogt{} block maps token features $x_t$ to a transition and value update.
The transition is produced by a skew-symmetric connection
\[
    A_t = \Phi_{\mathrm{conn}}(x_t)
          - \Phi_{\mathrm{conn}}(x_t)^\top,
\]
then mapped to $U_t$ either by a matrix exponential, Cayley transform, or an
identity ablation. Values are produced by a projection
\[
    V_t = \Phi_{\mathrm{val}}(x_t).
\]
The accumulated state is normalized and projected back to the residual stream.

\subsection{Gates}

Experiments show that pure transport is insufficient for overwrite-style state
tracking. We therefore study two optional gates:
\begin{align}
    U_t &\leftarrow g^{(U)}_t U_t, \\
    V_t &\leftarrow g^{(V)}_t V_t.
\end{align}
The transport gate controls forgetting, while the value gate controls selective
writing. Gates may be scalar or rank-wise. Scalar gates are simpler, but they
can only reset the entire state at once. Rank-wise gates allow different
subspaces to be retained or cleared independently, which is important for
multi-slot routing.

We also evaluate several transport choices. Cayley/Magnus transport is natural
for norm-preserving motion but can hurt overwrite memory. Identity transport is
a strong diagnostic: if the identity variant works while Cayley fails, the
problem is not the affine scan itself but the chosen connection dynamics.

\paragraph{Coupled write-forget gate.}
The strongest variant ties writing and forgetting:
\begin{align}
    z_t &= \sigma(\Phi_z(x_t)), \\
    U_t &\leftarrow (1-z_t) U_t, \\
    V_t &\leftarrow z_t V_t.
\end{align}
This preserves the affine scan law while imposing the correct overwrite bias:
non-write tokens keep state nearly unchanged, and write tokens both erase the
old value and insert the new one.

\paragraph{Slot-addressed gate input.}
For tracked multi-slot tasks, the gate must decide whether the current write
belongs to the tracked slot. We therefore test a minimal addressing input for
the write gate:
\[
    z_t = \sigma(\Phi_z([x_t, x_{t-1}, x_{\mathrm{track}}])).
\]
Here $x_{t-1}$ exposes the slot token immediately before a value token, and
$x_{\mathrm{track}}$ is the prefix slot selected by the query. This is still a
scan-compatible recurrent operator; the extra signal changes gate
parameterization, not the affine recurrence.

\section{Related Work}

\paragraph{Attention and position extrapolation.}
The Transformer replaces recurrence with attention over all previous tokens
\cite{vaswani2017attention}. This makes content-based retrieval easy, but long
context behavior depends strongly on the positional mechanism. Our experiments
therefore include both RoPE and NoPE Transformer baselines. RoPE is the more
standard small-decoder baseline, while NoPE is a stronger control for
position-agnostic synthetic state tasks.

\paragraph{Linear attention and recurrent views of attention.}
Linear attention rewrites attention-like computation as recurrent state updates
\cite{katharopoulos2020transformers}. This line of work shows that some
attention mechanisms can be interpreted as fast recurrent updates. \mogt{}
differs in that the recurrent state is explicitly matrix-valued and evolves
through affine transport pairs $(U_t,V_t)$ with a scan composition law.

\paragraph{Structured state-space and gated recurrence.}
S4 and related state-space models demonstrate that carefully parameterized
linear systems can model long sequences efficiently \cite{gu2021s4}. Mamba
introduces selective state spaces, making the recurrent dynamics input-dependent
\cite{gu2023mamba}. GateLoop similarly studies data-controlled linear
recurrence \cite{koo2023gateloop}. Our findings are consistent with this
direction: fixed or ungated transport is too rigid for overwrite memory, and
input-dependent gates are necessary.

\paragraph{Other recurrent Transformer alternatives.}
RetNet and RWKV also pursue recurrent or retention-based alternatives to
attention while preserving parallelizable training or efficient inference
\cite{sun2023retnet,peng2023rwkv}. Compared with these works, this draft is
much narrower: it does not claim language-modeling superiority, but isolates a
specific affine transport mechanism and its failure modes on controlled state
tracking tasks.

\section{Experimental Setup}

\subsection{Models}

We compare small matched models to isolate mechanism rather than scale. Unless
otherwise stated, synthetic experiments use 2 layers, $d_{\mathrm{model}}=128$,
rank $r=16$, BF16 on one NVIDIA L4, and roughly 0.54M parameters. Transformer
baselines use causal self-attention with either RoPE or no positional encoding
(NoPE). NoPE is included because state-tracking tasks can be position-agnostic,
and RoPE can be an artificially weak extrapolation baseline.

\subsection{Tasks}

\paragraph{Single-slot last-value tracking.}
Sequences contain a small number of write events. Each write sets a latent
state to a value token. At the final query token, the model must output the most
recent value. Dense supervision optionally asks the model to predict the current
state after each write.

\paragraph{Tracked multi-slot last-value tracking.}
The first two tokens specify a tracked slot. The sequence then contains writes
to multiple slots. At the query token, the model must output the latest value
written to the tracked slot. This tests selective routing rather than a single
global memory.

\paragraph{Language modeling sanity.}
We also run a short-budget WikiText-style language-modeling protocol. These
results are included as a boundary condition: current \mogt{} does not beat the
matched scratch Transformer on this protocol.

\section{Results}

\subsection{Single-Slot Extrapolation From Context 128}

Table~\ref{tab:ctx128} shows the first positive signal. Gated \mogt{} trained at
context 128 extrapolates better than a matched RoPE Transformer. The result is
not a broad attention win because NoPE Transformer is much stronger than RoPE on
this task.

\begin{table}[h]
\centering
\caption{Train context 128, three seeds. Accuracy is mean $\pm$ standard deviation.}
\label{tab:ctx128}
\begin{tabular}{lccc}
\toprule
Eval context & Gated \mogt{} & RoPE Transformer & Gap \\
\midrule
128  & 100.00 $\pm$ 0.00 & 99.48 $\pm$ 0.25 & +0.52 \\
256  & 98.26 $\pm$ 3.02  & 70.25 $\pm$ 13.68 & +28.01 \\
512  & 89.49 $\pm$ 15.11 & 22.88 $\pm$ 8.35 & +66.60 \\
1024 & 69.71 $\pm$ 18.69 & 10.51 $\pm$ 1.24 & +59.20 \\
\bottomrule
\end{tabular}
\end{table}

NoPE Transformer reaches 68.64\% $\pm$ 15.51\% at context 1024 under the same
2048-example evaluation protocol, narrowing the gap substantially. This makes
NoPE or other position-robust attention a mandatory baseline for future
claims.

\subsection{Train-Context-512 Dense Supervision}

Table~\ref{tab:ctx512dense} is the strongest scaled diagnostic. We train at
context 512 with dense state supervision. Identity dual-gated \mogt{} reaches
100\% at contexts 512 and 1024 across all seeds, but it decays faster than NoPE
Transformer at longer contexts. The coupled write-forget variant fixes this
specific failure mode on the single-slot task, staying at 100\% through context
8192 across all three seeds.

\begin{table}[h]
\centering
\caption{Train context 512 with dense state supervision, three seeds.}
\label{tab:ctx512dense}
\resizebox{\linewidth}{!}{%
\begin{tabular}{lcccc}
\toprule
Eval context & Coupled \mogt{} & Dual-gated \mogt{} & NoPE Transformer & Coupled gap \\
\midrule
512  & 100.00 $\pm$ 0.00 & 100.00 $\pm$ 0.00 & 97.40 $\pm$ 0.90 & +2.60 \\
1024 & 100.00 $\pm$ 0.00 & 100.00 $\pm$ 0.00 & 98.44 $\pm$ 1.56 & +1.56 \\
2048 & 100.00 $\pm$ 0.00 & 94.27 $\pm$ 8.61  & 95.31 $\pm$ 2.71 & +4.69 \\
4096 & 100.00 $\pm$ 0.00 & 70.83 $\pm$ 7.05  & 90.10 $\pm$ 5.49 & +9.90 \\
8192 & 100.00 $\pm$ 0.00 & 47.40 $\pm$ 8.88  & 78.65 $\pm$ 12.63 & +21.35 \\
\bottomrule
\end{tabular}
}
\end{table}

This table changes the story. The useful claim is not that \mogt{} dominates
attention. It is that the affine recurrent operator can learn long-context
state tracking once writing and forgetting are coupled. The remaining gap is
whether this mechanism scales to slot-addressed routing, language modeling, and
stronger recurrent baselines.

A seed-42 stress probe with fewer evaluation examples extends the same trained
setting to contexts 16k, 32k, and 65k, remaining at 100\% accuracy with 16
examples per context. We treat this as a capacity probe rather than a final
statistical result.

The learned gate is also interpretable. In the seed-42 coupled model, the first
block's mean write gate on a fresh training-context batch is 88.83\% on value
tokens, but only about $10^{-4}$\% on SET, QUERY, and filler tokens. Since the
same gate controls forgetting through $(1-z_t)$, this implements the intended
behavior: keep state on non-write tokens and overwrite on value tokens. The
second block is less selective, suggesting that the first recurrent layer
performs most of the explicit state update.

\subsection{Ablations}

Table~\ref{tab:ablation} summarizes seed-42 ablations. Cayley transport fails
on direct train512. Dense supervision helps. Identity transport plus dual gates
solves the train context and improves extrapolation. A value-gate-only variant
learns in-distribution but forgets faster at long context, suggesting that
explicit transport/forget gating is not just an optimization detail. Among the
non-coupled variants, NoPE Transformer has better far-context behavior and
lower memory in this implementation; the coupled variant closes the single-slot
quality gap but not the systems gap.

\begin{table}[h]
\centering
\caption{Seed-42 train512 stress probe.}
\label{tab:ablation}
\resizebox{\linewidth}{!}{%
\begin{tabular}{lccccc}
\toprule
Variant & 512 & 1024 & 2048 & 4096 & 8192 \\
\midrule
Gated \mogt{}, direct & 12.50 & 7.81 & 3.12 & 10.94 & 6.25 \\
\mogt{} Cayley dual-gate dense & 100.00 & 100.00 & 79.69 & 60.94 & 40.62 \\
\mogt{} identity dual-gate dense & 100.00 & 100.00 & 100.00 & 78.12 & 54.69 \\
\mogt{} identity value-gate-only dense & 87.50 & 92.19 & 64.06 & 48.44 & 25.00 \\
\mogt{} identity coupled write-forget dense & 100.00 & 100.00 & 100.00 & 100.00 & 100.00 \\
\mogt{} identity forget-ReLU dual-gate dense & 87.50 & 85.94 & 68.75 & 34.38 & 4.69 \\
NoPE Transformer dense & 98.44 & 100.00 & 96.88 & 89.06 & 82.81 \\
\bottomrule
\end{tabular}
}
\end{table}

\subsection{Tracked Multi-Slot State Routing}

Unconditioned multi-slot routing is weak: scalar, rank-wise, and coupled gates
without an explicit addressing signal stay near 50--67\% in the 2-slot probe.
Adding the slot-addressed gate input changes the result sharply.
Table~\ref{tab:multislot} shows the 3-seed 2-slot result.

\begin{table}[h]
\centering
\caption{Tracked 2-slot task, train context 512, dense supervision, 3000 steps, three seeds.}
\label{tab:multislot}
\resizebox{\linewidth}{!}{%
\begin{tabular}{lcccc}
\toprule
Model & 512 & 1024 & 2048 & 4096 \\
\midrule
Slot-addressed \mogt{} & 100.00 $\pm$ 0.00 & 100.00 $\pm$ 0.00 & 98.44 $\pm$ 1.56 & 86.46 $\pm$ 3.25 \\
NoPE Transformer & 45.31 $\pm$ 4.69 & 41.15 $\pm$ 3.93 & 34.90 $\pm$ 3.25 & 44.79 $\pm$ 9.42 \\
\bottomrule
\end{tabular}
}
\end{table}

The 4-slot result shows the same direction across three seeds: slot-addressed
\mogt{} reaches 100.00 $\pm$ 0.00, 100.00 $\pm$ 0.00, 100.00 $\pm$ 0.00, and
98.96 $\pm$ 0.90\% at contexts 512/1024/2048/4096, while NoPE Transformer
reaches 23.96 $\pm$ 3.61, 25.00 $\pm$ 4.13, 21.35 $\pm$ 4.77, and 24.48
$\pm$ 7.05\%.

Dense state supervision is not necessary for the easier 2-slot setting. With
final-query-only training, slot-addressed \mogt{} reaches 100.00 $\pm$ 0.00,
100.00 $\pm$ 0.00, 100.00 $\pm$ 0.00, and 96.35 $\pm$ 3.25\% at
512/1024/2048/4096 across three seeds, while NoPE Transformer reaches 42.19
$\pm$ 2.71, 40.62 $\pm$ 7.16, 47.92 $\pm$ 3.93, and 42.71 $\pm$ 21.67\%.
Direct 4-slot final-only training remains weak across three seeds: \mogt{}
reaches 44.79\% $\pm$ 17.21\% at 4096, compared with 25.00\% $\pm$ 5.63\%
for NoPE Transformer. However, a slot-count curriculum that trains on 2 active
slots and increases to all 4 active slots over the first 1500 steps changes the
result sharply (Table~\ref{tab:multislot4final}). This curriculum does not use
dense state labels; the loss is still only on the final query.

\begin{table}[h]
\centering
\caption{Tracked 4-slot task, final-query-only supervision with a 2-to-4-slot curriculum, train context 512, 3000 steps, three seeds.}
\label{tab:multislot4final}
\resizebox{\linewidth}{!}{%
\begin{tabular}{lcccc}
\toprule
Model & 512 & 1024 & 2048 & 4096 \\
\midrule
Slot-addressed \mogt{} & 100.00 $\pm$ 0.00 & 100.00 $\pm$ 0.00 & 99.48 $\pm$ 0.90 & 94.27 $\pm$ 2.39 \\
NoPE Transformer & 31.25 $\pm$ 5.41 & 29.69 $\pm$ 5.63 & 23.44 $\pm$ 5.63 & 21.35 $\pm$ 6.31 \\
HF-Mamba $d=192$ & 17.71 $\pm$ 10.97 & 20.83 $\pm$ 9.92 & 15.62 $\pm$ 5.63 & 19.79 $\pm$ 10.97 \\
\bottomrule
\end{tabular}
}
\end{table}

The same curriculum also extends to a 6-slot version across three seeds.
Slot-addressed \mogt{} reaches 100.00 $\pm$ 0.00, 100.00 $\pm$ 0.00,
98.96 $\pm$ 0.90, and 96.88 $\pm$ 2.71\% at contexts
512/1024/2048/4096, while NoPE Transformer reaches 20.83 $\pm$ 1.80,
22.92 $\pm$ 5.02, 21.35 $\pm$ 5.49, and 21.88 $\pm$ 8.27\%.
An 8-slot version remains positive but shows higher extrapolation variance:
slot-addressed \mogt{} reaches 100.00 $\pm$ 0.00, 100.00 $\pm$ 0.00,
97.40 $\pm$ 0.90, and 85.42 $\pm$ 11.93\% at contexts 512/1024/2048/4096,
while NoPE Transformer reaches 17.19 $\pm$ 2.71, 16.15 $\pm$ 6.51,
22.92 $\pm$ 5.02, and 10.94 $\pm$ 2.71\%. A small HF-Mamba baseline reaches
18.23 $\pm$ 3.61, 13.02 $\pm$ 1.80, 13.54 $\pm$ 7.86, and 18.23 $\pm$
1.80\%; a parameter-matched HF-Mamba baseline with $d=192$ reaches
17.71 $\pm$ 8.02, 13.54 $\pm$ 6.51, 19.27 $\pm$ 7.71, and 16.15 $\pm$
6.31\%. The 8-slot setting is therefore a useful scaling boundary: the
mechanism still works, but long-context extrapolation is no longer uniformly
saturated across seeds.
A seed-42 ablation without the active-slot curriculum reaches only 23.44\% at
4096, compared with 95.31\% for the curriculum run at the same seed. The
curriculum is therefore an optimization condition, not part of the recurrent
operator itself.
The same ablation is sharper at 8 slots: direct final-only training reaches
12.50\% at 4096, compared with 98.44\% for the curriculum run at the same
seed. Its gate diagnostic also fails to separate matched from unmatched value
tokens, suggesting that the curriculum is needed to discover the addressing
solution.

As recurrent controls, we added both a scratch 2-layer GRU and a small
HF-Mamba baseline. In a seed-42 500-step early-learning probe on the 6-slot
curriculum, the GRU reaches 6.25\% training accuracy at step 500 and 6.25\%
evaluation accuracy at context 4096. The HF-Mamba baselines are more useful SSM
controls and are included in the 8-slot comparison above, where both remain far
below \mogt{}.

A gate diagnostic confirms the intended mechanism. In a seed-42 2-slot run,
block 0 opens its mean rank-wise value gate to 32.64\% on value tokens whose
preceding slot matches the tracked prefix, but only 2.04\% on value tokens from
other slots and below 0.2\% on SET, slot, filler, and QUERY tokens. Block 1 is
less selective, again suggesting that the first recurrent layer performs the
explicit state-write operation.

\subsection{Language Modeling Boundary}

The current language-modeling results are negative for the broad replacement
claim. On a 32k-context, 200-step WikiText-style protocol, the scratch
Transformer baseline reaches loss 6.1009, while the \mogt{} baseline-v1
three-seed mean is 6.4067. A fair paper must present this boundary clearly.

We have also added an initial hybrid MOGT/Transformer causal LM path. In a
one-seed wiring pilot at context 8192 with two layers and
$d_{\mathrm{model}}=128$, an alternating hybrid reaches validation loss 10.4828
after 10 steps, compared with 10.5214 for the same-size Transformer. This is
not a language-modeling quality result; it only establishes that hybrid layers
can be trained under the same data and reporting protocol, motivating a
matched layer-ratio sweep.

A first four-layer ratio pilot gives the same kind of weak but useful signal:
with 0/25/50/75/100\% \mogt{} layers, validation loss after 5 steps is
10.6355/10.6159/10.6100/10.5838/10.5526. This monotone pattern is encouraging,
but the run uses one seed and one validation batch, so it remains a scheduling
signal rather than a result. Moreover, after zero-initializing the attention
output projection to match the identity-like start of \mogt{} residual
branches, the 0\% control improves to 10.6032 and the hybrid/all-\mogt{} points
cluster near 10.55. Thus this pilot mainly motivates larger controlled sweeps;
it does not yet identify an optimal layer ratio.
Extending only the 0/25/100\% zero-initialized controls to 50 steps gives
validation loss 9.5796/9.5618/9.6249, respectively. This makes pure \mogt{} a
negative control at the current budget and points instead toward low-ratio
hybrid insertion.
Across three seeds at 50 steps, 0/25/50\% \mogt{} layers obtain mean validation
loss 9.5865/9.5784/9.6034 with standard deviations 0.0091/0.0163/0.0167. This
keeps the 25\% insertion point alive, but the effect is small and
variance-overlapping.
Across three seeds, a 200-step follow-up with four validation batches
strengthens the candidate: 0\% and 25\% \mogt{} reach validation loss
7.6069 $\pm$ 0.0091 and 7.4897 $\pm$ 0.0107, respectively. This is the first
clear small-LM positive signal for low-ratio insertion, but it remains a
WikiText-103 pilot rather than a broad language-modeling result.
A same-budget single-layer position ablation suggests the insertion point
matters: attention-only gives 7.6069 $\pm$ 0.0091, while placing the single
\mogt{} layer at layer 1/2/3 gives 7.4897 $\pm$ 0.0107,
7.4675 $\pm$ 0.0110, and 7.4539 $\pm$ 0.0092, respectively. Layer 0 currently
has only seed 42 and gives 7.5189. This points toward late-layer insertion as
the next scale-up target, not as a final architecture rule.
The first such scale-up keeps the signal alive: with 500 optimizer steps and
eight validation batches, attention-only reaches 6.8137 $\pm$ 0.0139, while a
single layer-3 \mogt{} insertion reaches 6.7241 $\pm$ 0.0132. This is stronger
than the 50/200-step probes, but remains a small WikiText-103 experiment.
A 1000-step continuation in the same setting gives 6.3868 $\pm$ 0.0098 for
attention-only and 6.3357 $\pm$ 0.0128 for layer-3 \mogt{}. The gap narrows but
persists, making larger-model scale-up the next meaningful falsification test.
The first width scale probe is also positive: at $d_{\mathrm{model}}=192$,
four layers, 500 steps, and eight validation batches, attention-only reaches
6.5024 $\pm$ 0.0141 while layer-3 \mogt{} reaches 6.4544 $\pm$ 0.0134. The
hybrid has slightly fewer parameters in this setting, so the effect is not
explained by a larger parameter count.
However, the same width at 1000 steps is a boundary: attention-only reaches
6.0997 $\pm$ 0.0126, layer-2 \mogt{} nearly matches it at
6.1049 $\pm$ 0.0147, and layer-3 \mogt{} trails at
6.1337 $\pm$ 0.0162. A two-\mogt{} middle hybrid using layers 1 and 2 reaches
6.1416 $\pm$ 0.0180, so simply increasing the \mogt{} layer count does not fix
the boundary. Thus the present evidence supports a promising hybrid
optimization signal, not a durable language-modeling scaling law.
A learning-rate fairness check further weakens the simple hybrid claim:
at $d_{\mathrm{model}}=192$, 1000 steps, and learning rate $5\cdot10^{-4}$,
attention-only reaches 5.8753 $\pm$ 0.0108, while layer-2 \mogt{} reaches
5.9058 $\pm$ 0.0430. The higher learning rate improves both models, but the
tuned attention control remains stronger.
A residual-scale diagnostic suggests the boundary is partly a fusion problem:
in the same seed-42 layer-2 setting, \mogt{} readout residual scales
0.25/0.5/0.75/1.0 give validation losses
5.9199/5.9110/5.9174/5.9545, while the same-seed attention-only control
remains better at 5.8877. A three-seed fixed-scale confirmation at scale 0.5
nearly closes the gap: attention-only reaches 5.8753 $\pm$ 0.0108, while the
layer-2 \mogt{} hybrid reaches 5.8775 $\pm$ 0.0293. The scaled hybrid beats
attention on two paired seeds, but its mean is still 0.0022 worse, so this is a
stabilization result rather than a language-modeling win.
The naive learned scalar gate is negative in the same setting: initialized at
0.5, it reaches validation loss 5.9752, worse than the fixed 0.5 residual
scale. A simple bounded residual schedule is also negative: warming the readout
scale from 0.25 to 0.5 over 250 steps reaches 5.9497. Future residual mixing
therefore needs richer dynamics or auxiliary supervision rather than an
unconstrained scalar gate or a naive linear scale warmup.
The strongest follow-up is optimizer partitioning. With fixed residual scale
0.5 but a 0.5 learning-rate multiplier on MOGT block parameters, the seed-42
layer-2 hybrid reaches validation loss 5.8601, beating the same-seed
attention-only control at 5.8877. The three-seed confirmation is also positive:
attention-only reaches 5.8753 $\pm$ 0.0108, while the layer-2 hybrid with
\mogt{} learning-rate multiplier 0.5 reaches 5.8511 $\pm$ 0.0117, winning on
all paired seeds. This is the first matched small-budget language-modeling
win, but it still requires larger-scale confirmation.
The first width migration diagnostic is neutral: at $d_{\mathrm{model}}=256$,
seed 42, and the same budget, attention-only reaches 5.7550 while the
corresponding layer-2 \mogt{} hybrid reaches 5.7558. Thus the recipe does not
collapse at the next width, but it also does not yet provide a width-scaling
win.

\begin{table}[h]
\centering
\caption{WikiText-103 wiring pilot, context 8192, two layers, $d_{\mathrm{model}}=128$, seed 42, 10 optimizer steps.}
\label{tab:lm-pilot}
\begin{tabular}{lcc}
\toprule
Model & Best val loss & PPL \\
\midrule
Transformer & 10.5214 & 37099.21 \\
Alternating hybrid & 10.4828 & 35694.57 \\
\bottomrule
\end{tabular}
\end{table}

\subsection{Systems Boundary}

We also run two bounded systems probes on one NVIDIA L4. These are not
end-to-end throughput claims. The first times only the core affine scan and the
core attention operation; it excludes connection/value projection, the
connection map, FFN, normalization, LM head, loss, and backward pass. The
second times backbone hidden-state forward only: embeddings, sequence blocks,
and final normalization, without LM head, loss, backward pass, optimizer, or
KV-cache decode behavior.

\begin{table}[h]
\centering
\caption{Core operator timing, $d_{\mathrm{model}}=768$, rank 16, batch size 1. Lower is better.}
\label{tab:systems-core}
\begin{tabular}{lccc}
\toprule
Length & Affine Triton hybrid & Attention core & Attention / affine \\
\midrule
8192  & 3.12 ms & 1.48 ms  & 0.47$\times$ \\
16384 & 5.36 ms & 6.39 ms  & 1.19$\times$ \\
32768 & 5.48 ms & 27.32 ms & 4.98$\times$ \\
\bottomrule
\end{tabular}
\end{table}

\begin{table}[h]
\centering
\caption{Backbone hidden-state forward timing, two layers, $d_{\mathrm{model}}=768$, batch size 1. Lower is better.}
\label{tab:systems-backbone}
\begin{tabular}{lccc}
\toprule
Length & Identity coupled \mogt{} & NoPE Transformer & Transformer / \mogt{} \\
\midrule
8192  & 14.07 ms & 10.81 ms & 0.77$\times$ \\
16384 & 37.88 ms & 35.29 ms & 0.93$\times$ \\
32768 & 77.56 ms & 97.05 ms & 1.25$\times$ \\
\bottomrule
\end{tabular}
\end{table}

The system story is therefore promising but still immature. The affine scan
core has a strong long-context signal, and the small backbone forward probe
crosses over at 32k. However, the current \mogt{} backbone uses more peak
memory in this probe and the measurement does not cover training throughput or
autoregressive decoding.

\section{Reproducibility}

All experiments in this draft are produced by scripts in the repository. The
main synthetic scripts are:
\begin{itemize}
    \item \texttt{benchmark\_synthetic\_last\_value.py}
    \item \texttt{benchmark\_synthetic\_multislot.py}
    \item \texttt{train\_budget\_hybrid.py}
    \item \texttt{summarize\_synthetic\_last\_value.py}
    \item \texttt{summarize\_standard\_multislot.py}
    \item \texttt{validate\_experiment\_reports.py}
\end{itemize}

The primary result artifacts are:
\begin{itemize}
    \item \texttt{benchmark\_runs/synthetic\_last\_value\_summary\_20260503.md}
    \item \texttt{benchmark\_runs/synthetic\_multislot\_standard\_summary\_20260504.md}
    \item \texttt{benchmark\_runs/standard\_report\_index.md}
    \item \texttt{benchmark\_runs/budget\_matched\_baseline\_summary.md}
    \item \texttt{benchmark\_runs/hybrid\_lm\_ratio\_sweep\_summary\_20260504.md}
    \item \texttt{benchmark\_runs/hybrid\_layer\_position\_summary\_20260504.md}
    \item \texttt{benchmark\_runs/hybrid\_layer\_position\_steps500\_summary\_20260504.md}
    \item \texttt{benchmark\_runs/hybrid\_layer\_position\_steps1000\_summary\_20260504.md}
    \item \texttt{benchmark\_runs/hybrid\_scale\_d192\_l4\_steps500\_summary\_20260504.md}
    \item \texttt{benchmark\_runs/hybrid\_scale\_d192\_l4\_steps1000\_summary\_20260504.md}
    \item \texttt{benchmark\_runs/hybrid\_scale\_d192\_l4\_steps1000\_lr5e4\_summary\_20260504.md}
    \item \texttt{benchmark\_runs/hybrid\_scale\_d192\_l4\_steps1000\_lr5e4\_mogtscale\_sweep\_seed42\_20260505.md}
    \item \texttt{benchmark\_runs/hybrid\_scale\_d192\_l4\_steps1000\_lr5e4\_mogtscale0p5\_summary\_20260505.md}
    \item \texttt{benchmark\_runs/hybrid\_scale\_d192\_l4\_steps1000\_lr5e4\_mogtgate0p5\_diagnostic\_seed42\_20260505.md}
    \item \texttt{benchmark\_runs/hybrid\_scale\_d192\_l4\_steps1000\_lr5e4\_mogtsched0p25to0p5s250\_diagnostic\_seed42\_20260505.md}
    \item \texttt{benchmark\_runs/hybrid\_scale\_d192\_l4\_steps1000\_lr5e4\_mogtscale0p5\_mogtlr0p5\_diagnostic\_seed42\_20260505.md}
    \item \texttt{benchmark\_runs/hybrid\_scale\_d192\_l4\_steps1000\_lr5e4\_mogtscale0p5\_mogtlr0p5\_summary\_20260505.md}
    \item \texttt{benchmark\_runs/hybrid\_scale\_d256\_l4\_steps1000\_lr5e4\_mogtscale0p5\_mogtlr0p5\_summary\_seed42\_20260505.md}
\end{itemize}

The code paths were checked with affine scan sanity tests and Triton hybrid
training regression tests. PDF compilation was not verified in this environment
because \texttt{pdflatex} was not installed.

\section{Discussion}

\paragraph{Why identity transport helps.}
Overwrite-style memory requires preserving a value indefinitely until a new
write arrives. Cayley/Magnus transport rotates state, which is useful for some
dynamical systems but can make simple persistence difficult. Identity transport
is the correct inductive bias for this family of tasks.

\paragraph{Why NoPE Transformer is strong.}
The single-slot and multi-slot tasks are largely position-agnostic. A RoPE
Transformer trained at short context can fail due to positional extrapolation,
whereas NoPE attention can attend by content without depending on absolute
position. This is why RoPE-only comparisons are insufficient.

\paragraph{What remains missing.}
The current \mogt{} implementation is not yet a top-tier model result. It lacks
non-synthetic evidence, the hardest final-query-only 4-slot result requires a
slot-count curriculum, and language modeling quality is behind a matched
scratch Transformer. The strongest contribution is mechanistic: it identifies
write-aligned forgetting and prefix-conditioned addressing as missing
ingredients for affine scan operators on recurrent state tracking.

\section{Limitations}

This draft has several important limitations.
\begin{itemize}
    \item The strongest positive results are synthetic. They do not imply
    language-modeling superiority.
    \item NoPE Transformer is a very strong baseline for these state-tracking
    tasks. The coupled write-forget variant beats it on single-slot dense
    tracking, but this does not imply general attention replacement.
    \item Some of the strongest results use dense state supervision. The
    strongest 4/6/8-slot final-only results remove dense labels but use an
    explicit slot-count curriculum, which is still not ordinary language-model
    training.
    \item The current Triton path is a hybrid implementation, not a final fused
    production kernel.
    \item Multi-slot routing is solved in tracked, prefix-conditioned 2-slot
    final-only training, 4-slot dense-supervision training, and 4/6/8-slot
    final-only training with a slot-count curriculum. Direct final-only
    learning without curriculum and non-synthetic tasks remain open.
    \item The experiments are small. A submission-quality result needs larger
    model scales, more seeds for every table, stronger recurrent baselines, and
    real language-modeling or hybrid-attention evidence.
    \item The scratch GRU control is only an early-learning probe. The HF-Mamba
    control is stronger, but a submission should still include more optimized
    recurrent or linear-attention baselines where practical.
\end{itemize}

\section{Conclusion}

Matrix-valued affine transport provides an associative scan-compatible
alternative to attention. In its ungated Cayley/Magnus form, it does not solve
the tasks studied here. With identity transport, dense state supervision, and
a coupled write-forget gate, it learns single-slot recurrent state tracking and
extrapolates from context 512 to 8192 across three seeds. With a
prefix-conditioned slot-addressed gate, it also learns tracked multi-slot
routing in controlled 2-slot, 4-slot, 6-slot, and 8-slot probes, including
4/6/8-slot final-query-only training when slots are introduced by curriculum.
These results define the next target: transfer the mechanism to harder routing,
language modeling, or hybrid attention settings while retaining the scan
structure.

\bibliographystyle{plain}
\bibliography{references}

\end{document}
